\begin{abstract}
\color{red} 
The proliferation of diffusion-based generative models has intensified the need for reliable methods to trace and verify the provenance of AI-generated images. Existing watermarking approaches for diffusion models typically rely on isotropic frequency patterns and handcrafted detection metrics, which suffer from directional vulnerability and threshold instability across diverse image transformations. We propose SpiderMark, a hybrid watermarking framework that couples algorithmic frequency-domain injection with learning-based probabilistic verification, requiring no modification to pretrained diffusion model weights. SpiderMark embeds a structured spider-web pattern, comprising radial spokes and concentric polygonal rings, into the Fourier representation of the initial latent noise, and trains a lightweight ResNet-18 verifier augmented with auxiliary similarity cues (PSNR and L1) to produce calibrated confidence scores for watermark detection. Extensive experiments across nine image transformations demonstrate that SpiderMark consistently outperforms Tree-Ring, DFT-based, and DWT DCT baselines in both accuracy and AUROC, while generalising effectively to the unseen COCO dataset without retraining. Semantic evaluation via CLIP similarity and FID scores confirms that robustness is achieved without compromising image quality or prompt fidelity. SpiderMark establishes that structured watermark geometry and learned probabilistic verification can be unified into a single, deployment-ready framework, bridging the gap between imperceptibility, robustness, and real-world verifiability for generative AI content attribution.
\end{abstract}